\documentclass[10pt,a4paper]{amsart}
\usepackage[margin=19mm]{geometry}
\usepackage{amsmath,amssymb,mathtools}
\usepackage[scr=rsfs,cal=euler]{mathalfa}
\usepackage{xfrac}
\usepackage[unicode,hidelinks,bookmarks=false]{hyperref}
\newcommand{\ie}{i.e.\ }
\newcommand{\cf}{cf.\ }
\newcommand{\resp}{respectively\ }
\newcommand{\Subsection}{\subsection}
\providecommand{\nobreakdash}{\nobreak\hbox{-}}
\setlength{\emergencystretch}{2em}
\multlinegap0pt
\usepackage{amssymb}
\usepackage{dsfont}
\usepackage{tikz}
\newtheorem{theorem}{Theorem}
\title{Connes spectral distances, quantum discord and coherence of qubits}
\author{Bing-Sheng Lin, Zi-Hao Xu, Ji-Hong Wang, Han-Liang Chen}
\date{Source: arXiv:2206.10527v3 (CC BY 4.0)}
\begin{document}
\maketitle

\noindent\emph{Source and licence.} Connes spectral distances, quantum discord and coherence of qubits. Version arXiv:2206.10527v3 (CC BY 4.0), distributed by arXiv under the Creative Commons Attribution 4.0 International licence, http://creativecommons.org/licenses/by/4.0/, as recorded in the arXiv licence metadata for this exact version and shown on the abstract page https://arxiv.org/abs/2206.10527v3. Full licence notice: \texttt{licenses/arxiv-2206.10527-CC-BY-4.0.txt}.\par\smallskip

\noindent\textit{Editorial scope.} Counted unit: the article's theorem drr (source0.tex lines 448--463). The article's complete original proof of it is delivered with it (source0.tex lines 465--493, one \texttt{proof} environment of 29 lines). No mathematics is written, repaired or re-derived: the statement and the proof are the article's own lines, in the article's own notation, and no retained line is substituted or re-flowed.  Retained uncounted as the source-local context the proof needs: lines 408--412; lines 436--441.  Nothing was omitted from any retained span, so the package carries no omission record.  No retained line contains a citation, so the wrapper carries no bibliography and no bibliography entry.  No retained line contains a figure environment, an image-inclusion command or drawing code, so no image is delivered and the package declares no asset dependency.  Editorial formatting only, in addition: an AMS \texttt{amsart} preamble that restates the article's own declarations and macros used by the retained lines, namely the environments \texttt{theorem} on separate counters, together with the standard packages those lines need.  The retained environments print in this document as Theorem 1 in order of appearance, whereas the article prints the same items with the numbering of its own full document.  The complete native source of the article as distributed in the arXiv e-print for this exact version is bundled unchanged as \texttt{sources/127236/source0.tex}.
\begin{equation}\label{sd1}
	d(\rho_1,\rho_2)
	=\frac{1}{2}\sqrt{\frac{\hbar}{2}}\frac{(\Delta x)^2+(\Delta y)^2+(\Delta z)^2}{|\Delta z|}
	=\frac{1}{2}\sqrt{\frac{\hbar}{2}}\frac{|\Delta \vec{r}|^2}{|\Delta z|}.
\end{equation}

\begin{equation}\label{sd2}
	d(\rho_1,\rho_2)
	=\frac{1}{2}|(s-t)\Delta z|
	+|w|\sqrt{(\Delta x)^2+(\Delta y)^2}
	=\sqrt{\frac{\hbar}{2}}\sqrt{(\Delta x)^2+(\Delta y)^2}.
\end{equation}

\begin{theorem}
	In the spectral triple $(\mathcal{A},\mathcal{H},\mathcal{D})$, the Connes spectral distance between one-qubit states $\rho_1$ and $\rho_2$ is
	\begin{equation}\label{drr}
		d(\rho_1, \rho_2)=
		\left\{
		\begin{array}{ll}
			\sqrt{\dfrac{\hbar}{2}}\sqrt{(\Delta x)^2+(\Delta y)^2}
			=r\sin\theta\sqrt{\dfrac{\hbar}{2}},
			& \qquad\hbox{$\dfrac{\pi}{4}\leqslant \theta \leqslant \dfrac{3\pi}{4}$;} \\[1em]
			\dfrac{1}{2}\sqrt{\dfrac{\hbar}{2}}\dfrac{|\Delta \vec{r}|^2}{|\Delta z|}
			=\dfrac{r}{2|\cos\theta|}\sqrt{\dfrac{\hbar}{2}},
			& \qquad\hbox{others.}
		\end{array}
		\right.
	\end{equation}
\end{theorem}

\begin{proof}
It is easy to see that, the vector $\Delta \vec{r}$ is in the ball with $0\leqslant |\Delta \vec{r}|= r \leqslant 2$, and
\begin{equation}
	\sin \theta=\frac{\sqrt{(\Delta x)^2+(\Delta y)^2}}{\sqrt{(\Delta x)^2+(\Delta y)^2+(\Delta z)^2}}
	=\frac{\sqrt{(\Delta x)^2+(\Delta y)^2}}{r},\qquad
	\cos \theta=\frac{\Delta z}{r}.
\end{equation}

When
$(\Delta x)^2+(\Delta y)^2\leqslant (\Delta z)^2$, there is
$0\leqslant \sin\theta \leqslant\dfrac{\sqrt{2}}{2}$,
and $0\leqslant \theta \leqslant \dfrac{\pi}{4}$ or $\dfrac{3\pi}{4}\leqslant \theta \leqslant \pi$. From Eq.~(\ref{sd1}), we have
\begin{equation}
	d(\rho_1,\rho_2)
	=\frac{1}{2}\sqrt{\frac{\hbar}{2}}\frac{|\Delta \vec{r}|^2}{|\Delta z|}
	=\dfrac{r}{2|\cos\theta|}\sqrt{\dfrac{\hbar}{2}}.
\end{equation}


When $(\Delta x)^2+(\Delta y)^2\geqslant (\Delta z)^2$,
there is
$\dfrac{\sqrt{2}}{2}\leqslant \sin\theta \leqslant 1$,
and $\dfrac{\pi}{4}\leqslant \theta \leqslant\dfrac{3\pi}{4}$. From Eq.~(\ref{sd2}), we have
\begin{equation}
	d(\rho_1,\rho_2)
	=\sqrt{\frac{\hbar}{2}}\sqrt{(\Delta x)^2+(\Delta y)^2}
	=r\sin\theta\sqrt{\dfrac{\hbar}{2}}.
\end{equation}
\end{proof}

\end{document}
